\documentclass[main]{subfiles}
\usepackage[sols]{workbook}
\numbering{ws2-9}

\begin{document}

\ifSubfilesClassLoaded{%
  \chaptermark{\chaptertwo}
}{}

\section{Alternating Series Test} \label{ws2-9}

\begin{framed}
\centerline{\textbf{Lesson Goals}}

You should be able to:
\begin{itemize}
\item State the Alternating Series Test for Convergence, and explain the intuition behind it
\item Determine if a given series is an alternating series, and use the Alternating Series Test to determine if it converges
\item Find an upper bound on the error from a partial sum approximation of a convergent alternating series
\item Determine how many terms will ensure that a partial sum approximation of a convergent alternating series is within a certain margin of error
\item Test a series for absolute convergence to determine if it converges

\end{itemize}
\end{framed}

\begin{enumerate}

\item   \label{alt-harmonic}

  \begin{problem}[3.1in]
    Draw the first six partial sums of $\displaystyle \sum_{k=1}^{\infty} (-1)^{k+1} \frac{1}{k} $ on a number 
    line.  Does it look like the sum converges?
  \end{problem}

  \begin{solution}
    The partial sum diagram below
    suggests that the series converges, which actually is the case.  (Convergence
    follows from the Alternating Series Test.)
    \begin{center}
      \includegraphics[width=4in]{ws/images/alternating-harmonic}
    \end{center}
  \end{solution}

%  \item \begin{problem}[0.8in]
%    Can you relate the sum in part (a) to the Taylor series for $\mathrm{ln}(1+x)$ centered at 
%    $0$, evaluated at a specific value of $x$?
%  \end{problem}
  
%  \begin{solution}
%    The series is what you get if you plug $x=1$ into the Taylor Series for $\mathrm{ln}(1+x)$ 
%    centered at $x=0$.
%  \end{solution}



\begin{problemonly}
\begin{framed}
\centerline{\textbf{The Alternating Series Test}}

A series of the form $\displaystyle \sum_{k = 1}^{\infty} (-1)^{k} |b_{k}|$ or
$\displaystyle \sum_{k = 1}^{\infty} (-1)^{k+1} |b_{k}|$ where $|b_k| > 0$ for each $k$ is called an \uline{alternating series}.  If the following two conditions are met, then the alternating series converges:
\begin{itemize}
  \item $|b_{k+1}| \leq |b_{k}|$ for each $k = 1, 2, 3, \dots$
  \item $\displaystyle \lim_{k \to \infty} |b_{k}| = 0.$
\end{itemize}  
\end{framed}
\end{problemonly}

\pspace{\newpage}


%%%%%%%%%%%%%%%%%%%%%%%%%%%%%%%%%%%%%%%%%%%%%%%%%%%%%%%%%%%%%%%%%%%%%%%%%%%%%%%

\item     \label{test-alt-series}
\begin{problem}
Determine whether each of the series below converges or diverges.  Explain which
test for convergence or divergence you are using, and why the test applies.
\end{problem}

\begin{enumerate}
  \item  
  \begin{problem}[\vfill] 
     $\displaystyle\sum_{k=3}^\infty\frac{(-1)^k}{100\sqrt[3]{k}}$
  \end{problem}

  \begin{solution}
  The series does alternate, so let's see whether the Alternating Series Test
  applies.  Letting $\displaystyle b_k = \frac{1}{100 \sqrt[3]{k}}$, the series
  has the form $\displaystyle \sum_{k = 3}^{\infty} (-1)^{k} b_{k}$.  The first
  condition of the Alternating Series Test is met because
  $$
    b_{k+1} \; = \; \frac{1}{100 \sqrt[3]{k+1}} < \frac{1}{100 \sqrt[3]{k}} \; = \; b_{k}
  $$
  for each $k$.  The second condition of the Alternating Series Test is also
  met because
  $$
    \lim_{k \to \infty} b_{k} \; = \; \lim_{k \to \infty} \frac{1}{100 \sqrt[3]{k}} \; = \; 0.
  $$
  It follows that this series converges, by the Alternating Series Test.
  \end{solution}


  \item  
  \begin{problem}[\vfill]
     $\displaystyle\sum_{k=1}^\infty (-1)^{k+1} \frac{k}{2k + 1}$
  \end{problem}

  \begin{solution}
    Again, the series alternates, so let's see whether the Alternating Series Test
    applies.  Let $\displaystyle b_{k} = \frac{k}{2k + 1}$, and note that the
    series has the form $\displaystyle \sum_{k = 1}^{\infty} (-1)^{k+1} b_{k}$.
    But this time,
    $$
      \lim_{k \to \infty} b_{k} \; = \; \lim_{k \to \infty} \frac{k}{2k + 1} \; = \; \frac{1}{2}.
    $$
    This means that the Alternating Series Test does not apply---it is inconclusive.
    However, there is good news---the fact that the limit we just took is not equal
    to zero means that this series definitely diverges according to the $N$th term
    test for divergence.
  \end{solution}


\end{enumerate}

%%%%%%%%%%%%%%%%%%%%%%%%%%%%%%%%%%%%%%%%%%%%%%%%%%%%%%%%%%%%%%%%%%%

\begin{problemonly}
\begin{framed}
\textbf{Definition.}  A series $\sum a_k$ is called \uline{absolutely
convergent} if the series of absolute values $\sum |a_k|$ is
convergent.

\textbf{Theorem.}  If a series $\sum a_k$ is absolutely convergent,
then it is convergent.

The opposite (converse) of the theorem is not true: there are
convergent series which are not absolutely convergent.  These series
are called \uline{conditionally convergent}.
\end{framed}
\end{problemonly}

%%%%%%%%%%%%%%%%%%%%%%%%%%%%%%%%%%%%%%%%%%%%%%%%%%%%%%%%%%%%%%%%%%%%%%%%%%%%%%%


\item 
\begin{problem}[1in]
  Give an example of a series that converges but does not converge absolutely.
\end{problem}

\begin{solution} 
  The alternating harmonic series $\displaystyle \sum_{k=1}^{\infty} (-1)^{k+1}\frac{1}{k}$ is 
  convergent (by the Alternating Series Test, since $\displaystyle \frac{1}{k}$ is decreasing 
  and since $\displaystyle\lim_{k\rightarrow \infty} \frac{1}{k} = 0$).  But if we sum the 
  absolute values of the terms in the alternating harmonic series, we 
  get the harmonic series, which diverges.  Thus, the alternating harmonic series is an 
  example of a series that is \fbox{conditionally convergent}.
  
  Another example of a conditionally convergent series is in Problem~\ref{abs-conv}(b).
\end{solution}

\pspace{\newpage}


\item \label{abs-conv}

\note{In Part (a), I'll point out why it's advantageous to start by testing
for absolute convergence, and that there's never a need to invoke the
Alternating Series Test (even though it does apply).  Note that Part (b)
revisits the series from Problem~\ref{test-alt-series}(a) above.}


\begin{problem}
Determine whether each of the series below diverges, converges absolutely, or converges 
conditionally. Write down the test you are using. 
\end{problem}

\begin{enumerate}
  \item  
  \begin{problem}[\vfill] 
     $\displaystyle \sum_{k=1}^\infty (-1)^k\frac{4}{k^3}$
  \end{problem}

  \begin{solution}
  We will test for absolute convergence first because, if the series does
  converge absolutely, we'll be done (because absolute convergence implies
  convergence).  Summing the absolute values of the terms in the series,
  we obtain
  \[\sum_{k=1}^\infty \left|(-1)^k\frac{4}{k^3}\right| = \sum_{k=1}^\infty \frac{4}{k^3}.\]
  This is $4$ times a $p$-series with $p = 3 > 1$, which converges.  Therefore,
  the series $\displaystyle \sum_{k=1}^\infty (-1)^k\frac{4}{k^3}$  
  \fbox{converges absolutely}.  Note:  You could also use the Alternating
  Series Test to show that the series converges, but this would have been
  a waste of effort---it would have told you nothing about whether the
  series converges absolutely.
  \end{solution}



  \item  
  \begin{problem}[\vfill] 
     $\displaystyle\sum_{k=3}^\infty\frac{(-1)^k}{100\sqrt[3]{k}}$
  \end{problem}

  \begin{solution}
  This is the same series from Problem~\ref{test-alt-series}(a) above.
  We already know that the series converges by the Alternating Series Test,
  so let's test whether it converges absolutely.  Summing the absolute
  values of the terms, we obtain the series
  \[ \displaystyle\sum_{k=3}^\infty   \left|\frac{(-1)^k}{100\sqrt[3]{k}} \right| = 
  \displaystyle\sum_{k=3}^\infty\frac{1}{100\sqrt[3]{k}} = \frac{1}{100}  
  \sum_{k=3}^\infty\frac{1}{{k}^{1/3}}\]  This is a p-series with $p=1/3 \leq 1$, 
  which diverges.  Therefore, our original series does not converge absolutely.
  Since it converges but does not converge absolutely, the series is 
  \fbox{conditionally convergent}.
  \end{solution}


  \item  
  \begin{problem}[\vfill]
     $\displaystyle\sum_{k=1}^\infty \frac{- 3^k}{k^ {2/3}}$
  \end{problem}

  \begin{solution}
  This series doesn't alternate.  Moreover, $3^k \gg k^{2/3}$, and therefore 
  there is no way the terms are going to zero as $k \to \infty$.  By
  the Nth-term Test for Divergence, this series \fbox{diverges}.
  \end{solution}


  \item  
  \begin{problem}[\vfill] 
    $\displaystyle\sum_{k=1}^\infty \frac{(-3)^{k+ 20}}{4^ {k}}$
  \end{problem}

  \begin{solution}
  Notice right away that the series can be written as 
  \[\sum_{k=1}^\infty \frac{(-3)^{k+ 20}}{4^ {k}} = \sum_{k=1}^\infty 3^{20} 
  \left(\frac{-3}{4}\right)^{k}\]
  Which is a geometric series with common ratio $r=\frac{-3}{4}$.  If we sum
  the absolute values of the terms instead, we end up with a geometric series
  with common ratio $3/4$, which converges.  It follows that the series 
  $\displaystyle\sum_{k=1}^\infty \frac{(-3)^{k+ 20}}{4^ {k}}$ 
  \fbox{converges absolutely}.
  \end{solution}

\end{enumerate}





%%%%%%%%%%%%%%%%%%%%%%%%%%%%%%%%%%%%%%%%%%%%%%%%%%%%%%%%%%%%%%%%%%%

\pspace{\newpage}



\begin{problemonly}
\begin{framed}
\centerline{\textbf{Associated Error Bound for a Convergent Alternating Series:}}

Consider an alternating series
$$
\sum_{k=1}^\infty (-1)^{k+1} |a_k| = |a_1| - |a_2| + |a_3| - |a_4| + |a_5| - |a_6| + \cdots
$$
where $|a_k| > 0$ for each $k$.  If this series satisfies the conditions of the 
Alternating Series Test (and therefore converges to some limit $S$), then the 
error committed in using the partial sum
$$
  s_n \; = \;   |a_1| - |a_2| + |a_3| - |a_4| + \cdots + (-1)^{n+1} |a_n|
$$
to approximate $S$ is at most $|a_{n+1}|$, the magnitude of the first unused term.
That is, $$|s_n - S| \leq |a_{n+1}|$$
\end{framed}
\end{problemonly}


%%%%%%%%%%%%%%%%%%%%%%%%%%%%%%%%%%%%%%%%%%%%%%%%%%%%%%%%%%%%%%%%%%%



\item \label{sine1}
\note{Emphasize to students that we will once again assume that the sine function is equal to its Taylor series centered at $0$, which hasn't been fully verified so far.}

We have seen that the Taylor polynomials for $\sin(x)$ centered at 0 appear to give great approximations for the function when $x$ is close to 0. Assuming that these polynomial approximations actually get closer and closer to $\sin(x)$ as the degree increases, we would like to know how accurate these approximations actually are.

\begin{enumerate}

\item
\begin{problem}[\vfill] Show that the Taylor series for $\sin(x)$ centered at 0, given by $$\displaystyle \sum_{k=0}^{\infty} \frac{(-1)^k}{(2k+1)!}x^{2k+1}=x-\frac{1}{3!}x^3+\frac{1}{5!}x^5-\cdots$$ converges when $x=1$.   Do the same for $x=-1$.
\end{problem}

\begin{solution}
When $x=1$, the Taylor series for $\sin(1)$ centered at 0 is given by $\displaystyle \sum_{k=0}^{\infty} \frac{(-1)^k}{(2k+1)!}$. These terms are alternating, decreasing, and $\displaystyle\lim_{k\to\infty} \frac{1}{(2k+1)!}=0$. Therefore, by the Alternating Series Test, the series converges. We can show convergence similarly for $x=-1$.
\end{solution}

\item 
\begin{problem}[\vfill]
    Explain why the Taylor series for $\sin(x)$ converges for all $x$ in the interval $[-1,1]$.
\end{problem}

\begin{solution}
When $x=0$, the Taylor series clearly converges to 0. 

If $0<x<1$, the Taylor series will alternate in sign, with the first term being positive. And if $-1<x<0$, the Taylor series will also alternate in sign, but with the first term being negative. Thus, for all values of $x$, the Taylor series is an alternating series! Also, for all $x$ in $[-1,1]$, it is evident that the size of the terms $\displaystyle \left|\frac{(-1)^k\cdot x^{2k+1}}{(2k+1)!}\right|$ will decrease towards 0, since the numerator will not increase in size as $k\to\infty$, while the denominator will increase (quickly) to $\infty$ as $k\to\infty$. Thus, by the Alternating Series Test, the Taylor series for $\sin x$ converges for all $x$ in $[-1,1]$.
\end{solution}
\pspace{\newpage}
\item
\begin{problem}[\vfill] If you program a computer to compute $\sin(x)$ by using $\displaystyle \sin(x) \approx x - \frac{x^3}{3!} +  \frac{x^5}{5!}$ for all $x \in [-1, 1]$, what is an upper bound for the error involved in this approximation?
\end{problem}

\begin{solution}
We saw in the previous problem that for all the values of $x$ in this range, the terms of the series are alternating and decreasing to zero. So, we may use the alternating series error bound.  Thus, for any $x$, a bound on the error is 
\[  \left| \frac{x^7}{7!} \right| \]
What is the biggest this quantity can be on the interval $-1\leq x\leq 1$?  Evidently, the numerator is never bigger than when $x = \pm1$.  So a reasonable upper bound on the error is $\frac{1^7}{7!} =\frac{1}{5040}\leq \fbox{0.0002}$.  

This is actually really cool -- it means that if we use the  (easy to work with) degree $5$ Taylor polynomial for $\sin$ instead of the (hard to evaluate) function itself, we have less than an error of $0.0002$ on the interval $-1\leq x\leq 1$.
\end{solution}

\item 
\begin{problem}[\vfill]
If you want to approximate $\sin(1)$ with error less than $10^{-8}$, what \textit{degree} could you use for your Taylor polynomial to achieve this? 
\end{problem}

\begin{solution}
If $f(x) = \sin(x)$, then the Taylor series for $f$ centered at $0$ is given by
$$
  x - \frac{x^3}{3!} + \frac{x^5}{5!}-\frac{x^7}{7!} +\frac{x^9}{9!} - \frac{x^{11}}{11!} + \cdots
$$
Soon, we will have the tools needed to show that this series actually equals $\sin(x)$
for each real $x$.  Taking this for granted for now, this means
$$
  \sin(1) \; = \;  1 - \frac{1}{3!} + \frac{1}{5!} - \frac{1}{7!} + \frac{1}{9!} - \frac{1}{11!} + \cdots
$$
This is an alternating series which satisfies the conditions of the Alternating Series Test,
and therefore converges.  Moreover, given any two consecutive partial sums (say $s_{n}$ and $s_{n+1}$),
we know that whatever the series converges to must lie in between $s_{n}$ and $s_{n+1}$.  
Since $|s_{n+1} - s_{n}| = |a_{n+1}|$, this means that if we use $s_{n}$ to approximate
the sum of this alternating series, our error is at most as large as the magnitude of 
the first unused term, that is, $|a_{n+1}|$.

For instance, if we make the approximation
\begin{equation}
  \label{sine1-est}
  \sin 1 \approx 1 - \frac{1^3}{3!} + \frac{1^5}{5!}-\frac{1^7}{7!} 
   +\frac{1^9}{9!} - \frac{1^{11}}{11!} + \frac{1^{13}}{13!} - \frac{1^{15}}{15!}
\end{equation}
then our error is at most $\displaystyle \frac{1}{17!}$.  Moreover,
\begin{eqnarray*}
17! &=& (17 \cdot 16 \cdot 15 \cdot 14 \cdot 13 \cdot 12 \cdot 11 \cdot 10) \cdot 9!\\
    &>& (10 \cdot 10 \cdot 10 \cdot 10 \cdot 10 \cdot 10 \cdot 10 \cdot 10) \cdot 9!\\
    &=&  10^{8} \cdot 9! \; > \; 10^{8}.
\end{eqnarray*}
So the error in the approximation~(\ref{sine1-est}) above is less than 
\[\frac{1}{17!} < \frac{1}{10^8}\]
Or, if you happen to have a calculator handy, you can just compute: the error will be no 
greater than
\[\frac{1}{17!} \approx 2.81148 \times 10^{-15} < 10^{-8}\]

\fbox{So, we may use the degree $15$ Taylor polynomial.}  

\end{solution}









\end{enumerate}




%%%%%%%%%%%%%%%%%%%%%%%%%%%%%%%%%%%%%%%%%%%%%%%%%%%%%%%%%%%%%%%%%%%%%%%%%%%%%%%%%


%  NOTE FOR FALL 2023:  NOT SURE WHO WROTE THIS NEXT PROBLEM, BUT IT WASN'T WELL
%  THOUGHT-OUT AND I'M COMMENTING IT OUT.  NICE IDEA AND COULD BE ADAPTED, BUT
%  NEED TO GIVE STUDENTS GUIDANCE AS TO WHY THE MONOTONICITY CONDITION OF THE
%  ALT SERIES TEST IS MET FOR X OUTSIDE THE INTERVAL [-1,1].

%\item 
%\begin{problem}[2in]
%  If you program a computer to compute $\sin(x)$ by using $\displaystyle \sin(x) \approx x - 
%  \frac{x^3}{3!} +  \frac{x^5}{5!}$ for all $x \in [-2, 2]$, what is an upper bound for the 
%  error involved in this approximation?
%\end{problem}


%\begin{solution}
%For all the values of $x$ in this domain, the terms of the series are alternating and decreasing 
%to zero, so , as in the previous problem, we may use the alternating series error bound.  Thus, 
%for any $x$, a bound on the error is 
%\[  \left| \frac{x^7}{7!} \right| \]
%What is the biggest this quantity can be on the interval $-2\leq x\leq 2$?  Evidently, the 
%numerator is never bigger than when $x = \pm2$.  So a reasonable bound on the error is $2^7/7! 
%\leq 0.0254$  This is actually really cool -- it means that if we use the (super amazing, easy 
%to work with polynomial!) degree $5$ Taylor polynomial for $\sin$ instead of the (hard to 
%evaluate, pain in the butt) function itself, we have less than a $3\%$ error on the interval 
%$-2\leq x\leq 2$.

%\end{solution}


%%%%%%%%%%%%%%%%%%%%%%%%%%%%%%%%%%%%%%%%%%%%%%%%%%%%%%%%%%%%%%%%%%%%%%%%%%
\begin{comment}
\item 

\begin{problem}[2in]
Suppose you want to approximate $\cos(0.5)$ with error less than $0.0001$,
using a Taylor polynomial for the cosine function centered at $0$.
What degree polynomial can you take to guarantee the desired precision?
\end{problem}

\begin{solution}
We are assuming for now that the cosine function is equal to its Taylor series 
everywhere, i.e.,
\[
\cos x = 1 - \frac{x^2}{2!} + \frac{x^4}{4!} - \frac{x^6}{6!} + \cdots = \sum_{k=0}^{\infty} 
(-1)^k\frac{x^{2k}}{(2k)!}
\]
for each real $x$.

So we get that 
\[\cos (0.5) = 1 - \frac{(0.5)^2}{2!} + \frac{(0.5)^4}{4!} - \frac{(0.5)^6}{6!} + \cdots = 
\sum_{k=0}^{\infty} (-1)^k\frac{(0.5)^{2k}}{(2k)!}\]
Notice that the terms of the series are definitely decreasing to zero, since the numerator 
$(0.5)^{2k}$ decreases to $0$ and the denominator $(2k)!$ increases without bound. And the 
terms clearly alternate in sign.  So the conditions of the alternating series test are 
satisfied.  This tells is not only that the series converges, but that if we truncate the
series after finitely many terms, then the error is no greater 
than the size of the first unused term.  So, for example, 

\[\left| \,\, \cos(0.5) - \left[1 - \frac{(0.5)^2}{2!} + \frac{(0.5)^4}{4!} \right] \,\, 
\right| < \left| - \frac{(0.5)^6}{6!}\right| = \frac{(\frac{1}{2})^6}{6!}  = 
\frac{1}{(64)(720)}< \frac{1}{(50)(500)} = \frac{1}{25000} < \frac{1}{10000}\]


Suppose we change the question: how many terms do we need to find an approximation with error 
less than $1/10^{9}$?  We need to find a value of $N$ such that $\frac{1}{2^N\cdot 
N!}<\frac{1}{10^9}$, or alternately so that $2^N\cdot N!>10^9$.  Try coming up with a
``practical" choice of $N$ by hand.  You may want to look at the solution of
Problem~\ref{sine1} above for inspiration as to how to estimate factorials.

\end{solution}
\end{comment}



\end{enumerate}



\end{document}